\documentclass[11pt,a4paper]{article}

\usepackage[utf8]{inputenc}
\usepackage[T1]{fontenc}
\usepackage[english]{babel}
\usepackage[a4paper,margin=2.4cm]{geometry}

\usepackage{amsmath, amssymb, amsthm}
\usepackage{mathtools}
\usepackage{booktabs}
\usepackage{tabularx}
\usepackage{array}
\usepackage{caption}
\usepackage{float}
\usepackage{enumitem}
\usepackage{microtype}
\usepackage{xcolor}
\usepackage{listings}
\usepackage{fancyhdr}
\usepackage{titlesec}
\usepackage[hidelinks]{hyperref}

%  ---  Palette shared with the companion reports  ---
\definecolor{ink}{RGB}{26,30,36}
\definecolor{slate}{RGB}{92,101,112}
\definecolor{mist}{RGB}{233,235,238}
\definecolor{rulegrey}{RGB}{188,194,201}
\definecolor{steel}{RGB}{47,72,102}
\definecolor{steellight}{RGB}{124,146,172}
\definecolor{signal}{RGB}{160,32,45}

\hypersetup{
    colorlinks=true,
    linkcolor=steel,
    citecolor=steel,
    urlcolor=steel,
    pdftitle={Three Fleet Scenarios for ePlanSys},
    pdfauthor={Haniel Ulises},
}

\color{ink}

\setlength{\headheight}{22pt}
\pagestyle{fancy}
\fancyhf{}
\fancyhead[L]{\small\textcolor{slate}{\itshape Three Fleet Scenarios for ePlanSys}}
\fancyhead[R]{\small\textcolor{slate}{\thepage}}
\renewcommand{\headrulewidth}{0.4pt}
\renewcommand{\headrule}{\hbox to\headwidth{\color{rulegrey}\leaders\hrule height \headrulewidth\hfill}}

\titleformat{\section}{\normalfont\large\bfseries\color{ink}}{\thesection}{0.7em}{}
\titleformat{\subsection}{\normalfont\normalsize\bfseries\color{ink}}{\thesubsection}{0.7em}{}
\titleformat{\paragraph}[runin]{\normalfont\bfseries\color{slate}}{}{0em}{}[.]

\captionsetup{font=small, labelfont={bf,color=slate}, textfont=footnotesize}

\lstset{
  basicstyle=\ttfamily\footnotesize,
  keywordstyle=\color{steel},
  commentstyle=\color{slate}\itshape,
  frame=single,
  rulecolor=\color{rulegrey},
  breaklines=true,
  columns=fullflexible,
  showstringspaces=false,
}

\newcommand{\Kw}{\mathit{Kw}}

\title{\bfseries Three Fleet Scenarios for ePlanSys\\[2pt]
\large\normalfont\color{slate}EPDDL survey missions, grounded with \textsf{plank}
and solved with Aletheia}
\author{}
\date{}

\begin{document}
\maketitle
\thispagestyle{fancy}
\vspace{-2.2em}

\begin{abstract}
\noindent
ePlanSys shipped one worked example, the muddy children puzzle. It is the right
example for explaining what an epistemic goal \emph{is} and the wrong one for a
robotics user, who has no children and no mud. This note records three
robot-fleet scenarios written to replace it, at three sizes, together with what
Aletheia does on each. All three are checked into the ePlanSys tree as EPDDL
sources with matching action mappings; the two smaller ones are also regression
tests. The scenarios are deliberately the same three mechanisms at every size,
namely drive, look and say what you saw, so that what grows between them is
the search and not the modelling. Across the three, ground actions go from 8 to 48 and
expanded nodes from 17 to 5.03~million, while peak resident memory moves from
4.34~MB to 4.85~MB.
\end{abstract}

\section{Why a fleet}

The muddy children puzzle is a good teaching instance and a poor advertisement.
Nothing in it moves, nobody has a sensor, and the only action is a
simultaneous public announcement that no robot performs. A reader wanting to
know whether epistemic planning is worth the trouble on a real platform learns
nothing from it about what an action costs, what a sensor observes, or what a
radio does.

The scenarios below are built around a single situation that is common in fleet
robotics and inexpressible in PDDL:

\begin{quote}
A corridor may be blocked. Only a robot standing at the junction can see it.
The rest of the fleet has to route around it, so the rest of the fleet has to
know, and watching a robot look tells you that it looked, not what it saw.
\end{quote}

Three things follow, and they are the three mechanisms every scenario uses.
\emph{Driving} is ontic and public: it changes the world, and the fleet tracks
its members' positions. \emph{Looking} is semi-private sensing: the robot at
the junction learns the answer and the others learn only that a look happened.
\emph{Reporting} is a public announcement guarded by the speaker's knowledge,
so a robot can broadcast only what it actually found; because knowledge is
factive, hearing it settles the matter for everyone.

The goal is always of the form $\Kw_i(\text{blocked})$, knowing
\emph{whether} rather than knowing \emph{that}. This is not a stylistic choice. A goal
of $K_i(\text{blocked})$ is unachievable in the worlds where the corridor is
clear, so a plan for it would be unexecutable about half the time.

\section{The three scenarios}

\paragraph{Corridor} Two robots, one route. Either could drive out and look,
but looking is semi-private, so one robot sensing does not inform the other.
Both could go and look, at four actions; one can go, look and broadcast, at
three. The planner finds the second, which is to say it works out that talking
is cheaper than everybody going to see for themselves.

\paragraph{Depot} Three robots, two routes, and a depot rule: a robot
dispatched down a route stays on it. No robot can survey both routes, so the
fleet must split the work, and the robot that never moves ends up knowing both
answers purely by listening.

\paragraph{Survey} Four robots, three routes, same rules. Written to be big
rather than clever: it adds no mechanism, only size. Eight worlds the fleet
cannot tell apart to begin with, twelve goal conjuncts, and forty-eight ground
actions.

\begin{table}[H]
\centering
\small
\caption{The three shipped fleet scenarios, as \textsf{plank} grounds them and
as Aletheia solves them. Worlds is $|W|$ of the grounded initial model and all
of them are designated, so every world is a live possibility at planning time.
Goal conjuncts counts the $\mathit{Kw}$ atoms of the goal. Times are single
runs; see \S\ref{sec:method}.}
\label{tab:fleet-scenarios}
\begin{tabular}{@{}l r r r@{}}
\toprule
 & \textbf{Corridor} & \textbf{Depot} & \textbf{Survey} \\
\midrule
Robots                  & 2   & 3    & 4      \\
Routes                  & 1   & 2    & 3      \\
Atoms                   & 3   & 8    & 15     \\
Ground actions          & 8   & 24   & 48     \\
Worlds ($=$ designated) & 2   & 4    & 8      \\
Goal conjuncts          & 2   & 6    & 12     \\
\midrule
Solution depth          & 3   & 6    & 9      \\
Policy leaves           & 2   & 4    & 8      \\
Branches checked        & 4   & 12   & 28     \\
\midrule
Nodes expanded          & 17  & 1\,670 & 5\,029\,063 \\
Nodes generated         & 19  & 2\,010 & 5\,746\,557 \\
Search time (s)         & 0.00 & 0.01 & 58.82  \\
Peak RSS (MB)           & 4.34 & 4.47 & 4.85   \\
\bottomrule
\end{tabular}
\end{table}


\section{What the policies look like}

The corridor policy is the smallest thing that is recognisably a contingent
plan rather than a sequence:

\begin{lstlisting}
goto-junction_r1
  inspect_r1
    report-blocked_r1     <- if the camera saw debris
    report-clear_r1       <- if it did not
\end{lstlisting}

\noindent
The depot policy is the one worth reading twice, because the division of labour
is not something the domain states anywhere:

\begin{lstlisting}
goto-north_r1
  inspect-north_r1
    report-north-blocked_r1
      goto-south_r2
        inspect-south_r2
          report-south-blocked_r2
          report-south-clear_r2
    report-north-clear_r1
      goto-south_r2
        inspect-south_r2
          report-south-blocked_r2
          report-south-clear_r2
\end{lstlisting}

\noindent
\textsf{r1} takes the north route, \textsf{r2} the south, and \textsf{r3} never
moves. Nothing assigns those roles; the search arrives at them because the
depot rule makes any other allocation longer. \textsf{r3} finishes the mission
knowing the state of both corridors without having left the depot, which is the
entire argument for modelling communication epistemically rather than as an
ontic effect on a \texttt{told} predicate.

\section{Method and environment}
\label{sec:method}

Every instance was ground with \textsf{plank}
(\texttt{plank export -d $\ldots$ -p $\ldots$ -l intermediate.epddl}) and
solved with the Aletheia binary
(\texttt{epistemic\_planner -{}-task $\ldots$ -{}-plan $\ldots$}). Heuristic and
strategy were left to Aletheia's selection policy in all three cases; it chose
\emph{knowledge-spread} under the rule \texttt{kw-only-goal} and \emph{AO*}
under \texttt{sensing-small-designated} every time, which is the expected
reading of a $\Kw$-only goal over a small designated set. Every returned policy
was accepted by Aletheia's own validator.

Timings are single runs of \texttt{/usr/bin/time} on an otherwise idle machine:
Intel Core i7-10870H at 2.20\,GHz, 16 threads, 38\,GB RAM, Ubuntu 22.04.5,
kernel 6.8.0, GCC 11.4.0, Release build. They are not repeated measurements and
should be read as orders of magnitude, not as benchmarks. The sub-10\,ms
figures in particular are at the resolution floor of the tool.

\paragraph{Versions} ePlanSys \texttt{feature/epddl-front-end} at
\texttt{0050a87}; Aletheia at \texttt{d530d25}; \textsf{plank} at
\texttt{a682480}.

\section{Observations}

\paragraph{The search explodes and the representation does not} Expanded nodes
grow by roughly two orders of magnitude per size step: 17, then 1\,670, then
5.03~million, a factor of $98$ and then $3011$. Peak resident memory over the
same range grows by 12\,\%, from 4.34\,MB to 4.85\,MB. This is the property
Aletheia's design targets: models are three flat bit matrices and closed lists
store 128-bit fingerprints rather than states, so the cost of a large search is
time rather than memory. A planner storing states outright would not have
finished the third instance on this machine.

\paragraph{Depth is what hurts} The three instances differ in every parameter
at once, but the one that tracks the cost is solution depth, which AO* reaches
by iterative deepening: 3, 6, 9. Each survey trip adds three actions to the
depth and doubles the number of worlds, and the product of those two is what
the 3011$\times$ jump is made of.

\paragraph{The grounder produces multi-pointed models} Before these scenarios,
every task \textsf{plank} produced from the available EPDDL instances was
single-pointed, with one designated world, so sensing had exactly one possible
outcome and AO* returned a chain. The only branching fixture in the ePlanSys
tree was a file edited by hand. All three scenarios here ground to genuinely
multi-pointed models, because their \texttt{:init} theories leave the corridors
open rather than asserting a state for them. They are the first grounded
branching instances in that tree, and the first that exercise the executor's
branch node on a policy nobody hand-wrote.

\paragraph{Both routes to the planner agree} ePlanSys reaches Aletheia two
ways: \texttt{plansys2\_epistemic\_planner} links the search in, and
\texttt{plansys2\_aletheia\_plan\_solver} runs the binary as a subprocess and
reads its plan file back. On the corridor and depot instances the in-process
route reports node counts identical to the standalone binary, 17/19 and
1670/2010, and the subprocess route returns the same policy, translated
through the action mapping into PlanSys2 action expressions. The two are the
same planner and now have a shared test to say so.

\section{Reproducing}

\begin{lstlisting}
# ground
plank export -d robot-fleet-depot-domain.epddl \
             -p robot-fleet-depot-problem.epddl \
             -l intermediate.epddl -o /tmp/out

# solve
epistemic_planner --task /tmp/out/robot-fleet-depot-problem.json \
                  --plan /tmp/plan.json --explain
\end{lstlisting}

\noindent
Inside a built ePlanSys workspace the grounding step is implicit: the plan
solver is given the \texttt{.epddl} sources directly:

\begin{lstlisting}
planner:
  ros__parameters:
    EPISTEMIC:
      epddl_domain:  ".../robot-fleet-depot-domain.epddl"
      epddl_problem: ".../robot-fleet-depot-problem.epddl"
      action_mapping: ".../mappings/robot-fleet-depot.json"
\end{lstlisting}

\section{Limitations}

The timings are single runs on one machine and the three instances differ in
several parameters simultaneously, so nothing here isolates the effect of any
one of them; the scaling study in the companion report does that properly. The
survey instance takes about a minute, which is why it is shipped as an example
rather than as a regression test; the corridor and depot instances carry
that role. Nothing here has been executed on hardware: the policies are
validated by Aletheia and translated into PlanSys2 action expressions, but the
behaviour trees behind those expressions are stubs, so the claim is about
planning and not about a fleet that drove anywhere.

\end{document}
